%%%%%%%%%%%%%%%%%%%%%%%%%%%%%%%%%%%%%%%%%%%%%%%%%%%%%%%%%%%%%%%%%%%%%%%%%%%%%%%%
%%  EINLAGE 2: WISSENSCHAFTLICHE ARBEIT     Zwischenabgabe: 11.11.2026 (obligatorisch)
%%  Struktur: Beschreibung (1/4–1/2 S.) + Diskussion der Physik (1–3 S.) + Reflexion (1–2 S.)
%%  Paper-Eintrag in literatur.bib anlegen (Key z.B. Autor2020) und \PaperKey in einstellungen.tex setzen.
%%%%%%%%%%%%%%%%%%%%%%%%%%%%%%%%%%%%%%%%%%%%%%%%%%%%%%%%%%%%%%%%%%%%%%%%%%%%%%%%
% \PaperKey wird in einstellungen.tex gesetzt

\section{Einlage: Eine wissenschaftliche Arbeit zu einem Teilaspekt Ihres Schwerpunktgebietes}
\label{sec:paper}

\begin{hinweise}
Diese Einlage muss aus einer begutachteten Fachzeitschrift sein. Sie sollen mit dieser Einlage die
professionelle Literaturrecherche lernen. Bücher, Konferenzbeiträge, Trade Journals und ähnliches werden hier
nicht akzeptiert.
Dieser Artikel \textbf{muss} in Ihrer weiteren Arbeit verwendet werden (v.\,a. in Einlage 4: \enquote{Eine
Bezug zu Ihrem Paper ist Voraussetzung}).
Sie dürfen nicht den gleichen Artikel verwenden wie Ihre Partnerin / Ihr Partner.
Akzeptable wissenschaftliche Einlagen sind begutachtete wissenschaftliche Artikel aus wissenschaftlichen
Zeitschriften, die Sie über webofknowledge.com oder auch über Google Scholar finden. Erfüllt die Einlage die
Anforderung nicht, wird sie in der Bewertung nicht berücksichtigt. Bei Zweifeln über die Zulässigkeit der
beiden Einlagen mit wissenschaftlichen Arbeiten fragen Sie Prof.~O.~Marti.
\end{hinweise}

\subsection{Beschreibung}
\begin{hinweise}
Die Beschreibung gibt alle notwendigen Informationen an, die die Einlage beschreiben, also zum Beispiel den
Fundort, die Literaturstelle und eine kurze Inhaltsübersicht, aber keine Nacherzählung ($1/4$ Seite, max.
$1/2$ Seite!). Leitfaden: Art der Einlage, wie bin ich darauf aufmerksam geworden (Datenbank, Suchbegriffe,
Hinweis von \dots), wesentliche Aussagen in eigenen Worten.
\end{hinweise}
\begin{umfang}{Einlage 2 (Paper): Beschreibung}{0.25}{0.5}
\platzhalter{Vollständige Literaturstelle \cite{\PaperKey}; Fundort/Recherche; kurze Inhaltsübersicht;
Verweis auf Anhang~\ref{anh:paper}}
\end{umfang}

\subsection{Diskussion der Physik}
\begin{hinweise}
Beschreiben Sie hier die Physik dieser Einlage. Fassen Sie den Artikel zusammen. Arbeiten Sie besonders die
Physik, die für ihr Schwerpunktthema relavant ist, heraus.
\begin{itemize}\itemsep0pt
\item Was ist die Hauptaussage der wissenschaftlichen Arbeit?
\item Was sind die Voraussetzungen für den Inhalt des Artikels?
\item Wie hilft physikalisch die Physik dieser Einlage bei Ihrem Schwerpunktthema?
\item Wo gab es Schwierigkeiten beim Verständnis der Physik? Welche? Mit welcher zusätzlichen Literatur haben Sie das Problem gelöst
\item Haben Sie weitere Artikel oder Bücher konsultiert um das Paper zu verstehen? Welche?
\item Welche Aussagen verwenden Sie für Ihr Schwerpunktthema und wo? Warum?
\end{itemize}
\textbf{Die Diskussion der Physik sollte mindestens eine Seite und maximal 3 Seiten lang sein. Die
Längenangaben sind ohne Titel, Bilder, Tabellen und Tabellen- oder Bildlegenden und ohne übergrosse Abstände
zu Gleichungen zu verstehen.}
Gleichungen zur Argumentation verwenden, nummerieren, Symbole definieren; aus dem Paper übernommene Gleichungen
mit Gleichungsnummer zitieren: \texttt{\textbackslash cite[Gl.~(3)]\{\PaperKey\}}.
\end{hinweise}
\begin{umfang}{Einlage 2 (Paper): Diskussion der Physik}{1}{3}
\platzhalter{Hauptaussage, Voraussetzungen, relevante Physik für \Schwerpunkt}

% Beispiel für eine zitierte, nummerierte Gleichung (anpassen oder löschen):
% Nach \cite[Gl.~(1)]{\PaperKey} gilt
% \begin{equation}
%   p(r,\theta) = \dots
%   \label{eq:paper-hauptgleichung}
% \end{equation}
% wobei $p$ der Schalldruck, $r$ der Abstand ... ist. Aus Gl.~\eqref{eq:paper-hauptgleichung} folgt ...

\end{umfang}

\subsection{Reflexion}
\begin{hinweise}
In der Reflexion sollen Sie sich Gedanken machen über die Wirkung dieser Einlage und der Arbeit mit ihr auf
Ihren Lernfortschritt.
\begin{itemize}\itemsep0pt
\item In welchem Zusammenhang steht die Einlage mit Ihrem Schwerpunktthema?
\item Was haben Sie gelernt?
\item Wie hat Ihnen die Einlage beim Verständnis Ihres Projektes geholfen?
\item Wo gab es Schwierigkeiten bei der Beschäftigung mit der Einlage?
\item Hätte ein schon bestandenes Modul beim Verständnis helfen können? Welches? Warum oder warum nicht?
\item Was haben Sie durch diese Zusatzliteratur gelernt?
\item Was hätten Sie anders gemacht, wenn Sie die Einlage erstellt hätten?
\item Was sagen Ihre Kommilitoninnen oder Kommilitonen dazu, als Kommentar oder wenn Sie mit Ihnen über die Einlage diskutieren?
\item Haben Sie in dem Artikel der Einlage logische oder physikalische Fehler entdeckt?
\item Gab es bei der Beschäftigung mit der Einlage eine überraschende neue Erkenntnis?
\end{itemize}
\textbf{Die Reflexion für eine Einlage sollte mindestens eine Seite und maximal 2 Seiten lang sein. Die
Längenangaben sind ohne Titel, Bilder, Tabellen und Tabellen- oder Bildlegenden zu verstehen.}
\end{hinweise}
\begin{umfang}{Einlage 2 (Paper): Reflexion}{1}{2}
\platzhalter{Reflexion in Fließtext -- Leitfragen erkennbar beantworten}
\end{umfang}

%\subsection{Eigene Bewertungskriterien}
%...
